% !TEX root = ../report.tex
\chapter{Технологическая часть}
\section{Язык программирования и инструменты сборки}
Для реализации выбран язык C: он имеет статическую типизацию, поддерживает работу с массивами, указателями и рекурсивными функциями. Проверка выполнена компилятором GCC 16.1.1 с параметрами \texttt{-std=c11 -Wall -Wextra -Werror -Wpedantic}. Основная программа собирается с помощью CMake. Типы \texttt{size\_t} и \texttt{ssize\_t} используются для размера массива и результата поиска соответственно.

Подготовка отсортированной копии выполняется пузырьковой сортировкой вне функций поиска. Её затраты рассматриваются отдельно от числа сравнений одного поискового запроса. Элементы массива имеют тип \texttt{int}.

\section{Реализация алгоритмов}
Основные функции поиска приведены в листингах~\ref{lst:linear}--\ref{lst:recursive}. В алгоритме без ветвления и рекурсивной реализации вспомогательная функция \texttt{compare(a, b)} возвращает результат сравнения \texttt{a < b}.
\par\noindent\begin{minipage}{\linewidth}
\lstinputlisting[caption={Линейный поиск},label={lst:linear},firstline=1,lastline=12]{../src/linear_search.c}
\end{minipage}\par
\par\noindent\begin{minipage}{\linewidth}
\lstinputlisting[caption={Классический нерекурсивный бинарный поиск},label={lst:classic},firstline=1,lastline=18]{../src/binary_search_classic.c}
\end{minipage}\par
\par\noindent\begin{minipage}{\linewidth}
\lstinputlisting[caption={Бинарный поиск без ветвления},label={lst:branchless},firstline=1,lastline=30]{../src/binary_search_branchless.c}
\end{minipage}\par
\par\noindent\begin{minipage}{\linewidth}
\lstinputlisting[caption={Рекурсивный бинарный поиск},label={lst:recursive},firstline=1,lastline=24]{../src/binary_search_recursive.c}
\end{minipage}\par

\section{Учёт сравнений}
Счётчик учитывает каждое выполненное сравнение искомого значения с элементом массива. Каждое выполненное сравнение с помощью \texttt{==}, \texttt{<} или \texttt{>} увеличивает счётчик на единицу. Например, если сначала проверяется \texttt{a[M] < x}, а затем \texttt{a[M] > x}, учитываются два сравнения. В рекурсивном поиске складываются сравнения, выполненные во всех вызовах. Проверки границ и условий циклов, сортировка, ввод и вывод не учитываются.

\section{Тестирование}
В таблице~\ref{tab:tests} приведены тесты всех четырёх алгоритмов. Для линейного поиска используются исходные данные, для бинарных --- отсортированная копия.

\begin{table}[H]
\caption{Функциональные тесты алгоритмов поиска}\label{tab:tests}
\centering\tabfont
\begin{tabularx}{\textwidth}{|P{6mm}|Y|P{42mm}|Y|}
\hline
№ & Описание & Входные данные & Ожидаемый и полученный результат\\\hline
1 & Пустой массив & $a=[]$, $x=7$ & $-1$, сравнений нет\\\hline
2 & Единственный элемент найден & $a=[4]$, $x=4$ & Индекс 0\\\hline
3 & Значение меньше единственного & $a=[4]$, $x=3$ & $-1$\\\hline
4 & Значение больше единственного & $a=[4]$, $x=5$ & $-1$\\\hline
5 & Первый элемент & $a=[0,2,4,6,8]$, $x=0$ & Индекс 0\\\hline
6 & Средний элемент & $a=[0,2,4,6,8]$, $x=4$ & Индекс 2\\\hline
7 & Последний элемент & $a=[0,2,4,6,8]$, $x=8$ & Индекс 4\\\hline
8 & Отсутствие между элементами & $a=[0,2,4,6,8]$, $x=3$ & $-1$\\\hline
9 & Повторяющиеся значения & $a=[1,1,1,2]$, $x=1$ & Индекс из $\{0,1,2\}$\\\hline
10 & Границы типа \texttt{int} & \texttt{a = [INT\_MIN, 0, INT\_MAX]}, \texttt{x = INT\_MAX} & Индекс 2\\\hline
11 & Исходный порядок & $a=[8,0,4]$, $x=8$ & Линейный: 0; бинарные: 2\\\hline
\end{tabularx}
\end{table}

Все тесты пройдены успешно.


\section*{Вывод}
Реализации четырёх алгоритмов прошли функциональные проверки.
